\subsection{Reflexion und Diskussion}

HTTP/2 hat im Vergleich zu HTTP/1.1 einen Overhead bei der Initialisierung der Verbindung.
Nach diesem Overhead werden aber die Vorteile von HTTP/2 deutlich.
Durch HPACK sind die Header deutlich kleiner, was die Übertragung beschleunigt.
Außerdem ermöglicht HTTP/2 das Multiplexing von Anfragen, wodurch die Latenz speziell bei vielen parallel laufenden Anfragen deutlich reduziert wird.
Insgesamt zeigt sich, dass HTTP/2 deutlich schneller ist als HTTP/1.1, insbesondere bei vielen Anfragen oder großen Headern.
Die Ergebnisse bestätigen die theoretischen Vorteile von HTTP/2 und zeigen, dass die Implementierung von HTTP/2 in modernen Webbrowsern und Servern gerechtfertigt ist, um die Leistung und Benutzererfahrung im Web zu verbessern.
Moderne Webanwendungen profitieren von HTTP/2, da sie oft viele Ressourcen (Bilder, Skripte, Stylesheets) laden müssen, und HTTP/2 ermöglicht es, diese effizienter zu übertragen.
Ein Großteil der Seite kann bereits angezeigt werden, während die restlichen Ressourcen noch geladen werden, was die wahrgenommene Ladezeit verbessert.
Bei diesen Vorteilen fällt die anfängliche Verzögerung durch die Verbindungseinrichtung von HTTP/2 kaum ins Gewicht, zumal der Unterschied sogar im Direktvergleich mit HTTP/1.1 nicht für einen Nutzer spürbar ist.